%% kapa: exam=1 round=33 year=2022 period=1 form=A subject=economics subject_ko=경제학원론 seq=1 number=41
%% source: Q-Net 감정평가사 기출문제 (cst00309) artlSeq=5209997 file=raw/1008/5209997/2022년도 제33회 감정평가사 1차시험 1교시(민법,경제학,부동산학).pdf
%% conversion: kordoc 4.24.51 -> tools/gen1.py (auto); review=proofread-sonnet; flags=tikz
\begin{question}{41}
재화 X의 시장수요곡선($D$)과 시장공급곡선($S$)이 아래 그림과 같을 때, 균형가격(P$^*$)과 균형거래량(Q$^*$)은? (단, 시장수요곡선과 시장공급곡선은 선형이며, 시장공급곡선은 수평이다.)
\begin{center}
\begin{tikzpicture}[x=0.12mm,y=-0.12mm,font=\small,line width=0.7pt]
\draw[-{Stealth[length=1.5mm]}] (62,337) -- (62,10);
\draw[-{Stealth[length=1.5mm]}] (62,337) -- (388,337);
\draw (62,96) -- (302,337);
\draw (62,265) -- (375,265);
\node at (46,12) {\textbf{P}};
\node at (387,356) {\textbf{Q}};
\node at (176,183) {\textbf{\textit{D}}};
\node at (362,250) {\textbf{\textit{S}}};
\node at (33,96) {500};
\node at (33,265) {150};
\node at (44,337) {0};
\node at (302,355) {500};
\end{tikzpicture}
\end{center}
\begin{choices}
  \item P$^*$ = 150, Q$^*$ = 150
  \item P$^*$ = 150, Q$^*$ = 350
  \item P$^*$ = 150, Q$^*$ = 500
  \item P$^*$ = 350, Q$^*$ = 150
  \item P$^*$ = 500, Q$^*$ = 150
\end{choices}
\end{question}
\answer{2}
